% ---------------------------------------------------------------------------
% infra-invertibility — Invertibility interface: invertible = locally principal
% NOT in the thesis (used silently in §1.2 and §3.3). The definition below is
% the working definition frozen for the Lean development. Interface lemmas are
% sorry-stubs this pass; proofs land with WP-G.
% ---------------------------------------------------------------------------

Throughout, $d < 0$ is a discriminant ($d \equiv 0, 1 \pmod 4$), written
$d = D f^2$ with $D$ the fundamental discriminant of $K := \mathbb{Q}(\sqrt d)$
and $f \geq 1$ the conductor; $\sqrt d := f \sqrt D$. The maximal order is
$O_K = \mathbb{Z}[\omega]$, $\omega := \tfrac{D + \sqrt D}{2}$, and
$\mathcal{O} = \mathcal{O}_d := \mathbb{Z}[\tau]$, $\tau := \tfrac{d + \sqrt d}{2}$,
is the order of discriminant $d$: a free $\mathbb{Z}$-module with basis
$1, \tau$, with $g(x) = x^2 - dx + \tfrac{d^2 - d}{4}$ the minimal polynomial of
$\tau$. Conjugation $\alpha \mapsto \bar\alpha$ preserves $\mathcal{O}$
($\bar\tau = d - \tau$) and $N(\alpha) := \alpha \bar\alpha \in \mathbb{Z}$.
For a nonzero ideal $\mathfrak{a} \subseteq \mathcal{O}$ the index
$[\mathcal{O} : \mathfrak{a}] := \#(\mathcal{O}/\mathfrak{a})$ is finite and is
the ideal norm of the development. For an integral domain $R$ with fraction
field $F$ and a prime $\mathfrak{p} \subseteq R$ we realize the localization
inside $F$, as $R_{\mathfrak{p}} := \{ a/s : a \in R, \ s \in R \setminus \mathfrak{p} \}$,
and write $I R_{\mathfrak{p}}$ for the extension of an ideal $I \subseteq R$.
A maximal ideal $\mathfrak{m}$ of $\mathcal{O}$ has residue characteristic $q$
when $\mathfrak{m} \cap \mathbb{Z} = q\mathbb{Z}$ (such a $q$ always exists, by
the index layer); we say $\mathfrak{m}$ lies over $q$.

The thesis never isolates the following notion: its \S 1.2 speaks of ``locally
principal (invertible) ideals'' and the proof of Corollary 3.1.3 in \S 3.3
counts locally principal ideals while the statement counts invertible ones.
This section is the bridge, stated over a general integral domain where
possible.

\begin{definition}[Invertible ideal]\label{infra-invertibility-def}
  \lean{QuadraticOrder.IsInvertibleIdeal}
  \uses{notation}
  Let $R$ be an integral domain. An ideal $I \subseteq R$ is \emph{invertible}
  if there exist an ideal $J \subseteq R$ and an element $x \in R$, $x \neq 0$,
  such that
  \[ I J = x R . \]
  The Lean predicate \texttt{QuadraticOrder.IsInvertibleIdeal} is this
  definition for $R = \mathcal{O}_d$. No nonzeroness or finiteness hypothesis
  on $I$ is imposed; both are consequences
  (Proposition~\ref{infra-invertibility-arith}).
\end{definition}

\begin{remark}\label{infra-invertibility-def-remark}
  \uses{infra-invertibility-def}
  For a nonzero integral ideal $I$ this agrees with the textbook definition
  through fractional ideals (Neukirch I \S 12, Cox \S 7): if $I J = x\mathcal{O}$
  then $I \cdot (x^{-1} J) = \mathcal{O}$; conversely if $I \mathfrak{b} = \mathcal{O}$
  for a fractional ideal $\mathfrak{b}$, pick $0 \neq a \in I$ and set
  $J := a \mathfrak{b} \subseteq I\mathfrak{b} = \mathcal{O}$, so that $I J = a \mathcal{O}$.
  The inverse is unique and equals $(R :_F I)$. The definition was frozen in
  this integral form because Mathlib's fractional-ideal invertibility API is
  essentially Dedekind-gated, and $\mathcal{O}_d$ is not Dedekind for $f > 1$.
\end{remark}

\begin{proposition}[Arithmetic of invertible ideals]\label{infra-invertibility-arith}
  \uses{infra-invertibility-def}
  Let $R$ be an integral domain with fraction field $F$ and let
  $I, J, A, B \subseteq R$ be ideals.
  \begin{enumerate}
    \item[(a)] If $I$ is invertible then $I \neq 0$.
    \item[(b)] Every nonzero principal ideal is invertible; in particular $R$ is.
    \item[(c)] If $I$ and $J$ are invertible, so is $IJ$.
    \item[(d)] If $IJ$ is invertible, so are $I$ and $J$.
    \item[(e)] If $I$ is invertible and $IA = IB$ then $A = B$.
    \item[(f)] If $I$ is invertible then $I$ is finitely generated.
    \item[(g)] If $R$ is local and $I$ is invertible, then $I$ is principal,
      generated by an element of $I$.
  \end{enumerate}
\end{proposition}

\begin{proof}
  \uses{infra-invertibility-def}
  (a) If $I = 0$ then $IJ = 0 \neq xR$ for $x \neq 0$.
  (b) $(aR) \cdot R = aR$.
  (c) If $I J_1 = xR$ and $J J_2 = yR$ then
  $(IJ)(J_1 J_2) = xyR$ and $xy \neq 0$ in the domain $R$.
  (d) If $(IJ)L = xR$ then $I(JL) = xR$ and $J(IL) = xR$.
  (e) With $IJ = xR$: from $IA = IB$, multiplying by $J$ gives $xA = xB$;
  for $a \in A$ write $xa = xb$, $b \in B$, and cancel $x$ in the domain.
  Hence $A \subseteq B$, and symmetrically.

  For (f) and (g), let $IJ = xR$, $x \neq 0$. Since $x \in IJ$, write
  $x = \sum_{i=1}^n a_i c_i$ with $a_i \in I$, $c_i \in J$, and set
  $b_i := c_i / x \in F$. Then
  \[ \sum_{i=1}^n a_i b_i = 1
     \qquad\text{and}\qquad
     b_i I = \tfrac{1}{x} c_i I \subseteq \tfrac{1}{x} J I = R
     \quad (1 \leq i \leq n). \tag{$\ast$} \]
  (f) For $a \in I$: $a = \sum_i a_i (a b_i)$ with $a b_i \in R$ by $(\ast)$;
  so $I = (a_1, \dots, a_n)$.
  (g) Let $\mathfrak{m}$ be the maximal ideal. Each $a_i b_i$ lies in $R$ by
  $(\ast)$; were all of them in $\mathfrak{m}$, we would get
  $1 = \sum_i a_i b_i \in \mathfrak{m}$. So $u := a_{i_0} b_{i_0}$ is a unit
  for some $i_0$, and every $a \in I$ satisfies
  $a = a_{i_0} \cdot \bigl( u^{-1} (b_{i_0} a) \bigr) \in a_{i_0} R$
  since $b_{i_0} a \in R$. Hence $I = a_{i_0} R$ with $a_{i_0} \in I$.
\end{proof}

\begin{theorem}[Invertible iff locally principal]\label{infra-invertibility}
  \lean{QuadraticOrder.isInvertibleIdeal_iff_isLocallyPrincipal}
  \uses{infra-invertibility-def, notation, infra-index}
  Let $R$ be an integral domain and $I \subseteq R$ a nonzero finitely
  generated ideal. The following are equivalent:
  \begin{enumerate}
    \item[(i)] $I$ is invertible;
    \item[(ii)] $I R_{\mathfrak{p}}$ is principal for every prime ideal $\mathfrak{p}$;
    \item[(iii)] $I R_{\mathfrak{m}}$ is principal for every maximal ideal $\mathfrak{m}$;
    \item[(iv)] $I R_{\mathfrak{m}}$ is principal for every maximal ideal
      $\mathfrak{m} \supseteq I$;
    \item[(v)] $I R_{\mathfrak{m}}$ is an invertible ideal of $R_{\mathfrak{m}}$
      for every maximal ideal $\mathfrak{m}$.
  \end{enumerate}
  The implications (i) $\Rightarrow$ (ii), (iii), (iv), (v) hold without finite
  generation. Moreover, when these conditions hold, then with
  $I^{-1} := (R :_F I) = \{ \beta \in F : \beta I \subseteq R \}$ one has
  $I \, I^{-1} = R$, and for every nonzero $a \in I$ the ideal
  $J := a I^{-1} \subseteq R$ satisfies $I J = a R$.

  In particular, invertibility is a local property, and since $\mathcal{O}_d$
  is a Noetherian domain the equivalence applies to every nonzero ideal of
  $\mathcal{O}_d$; the Lean declaration is this instance.
\end{theorem}

\begin{proof}
  \uses{infra-invertibility-def, infra-invertibility-arith, infra-index}
  We use the standard localization facts, all inside $F$: elements of
  $I R_{\mathfrak{p}}$ have the form $a/s$ with $a \in I$, $s \notin \mathfrak{p}$
  (clear denominators); $\mathfrak{p} R_{\mathfrak{p}} \cap R = \mathfrak{p}$;
  $R_{\mathfrak{p}}$ is local with maximal ideal $\mathfrak{p} R_{\mathfrak{p}}$;
  extension of ideals is multiplicative,
  $(IJ) R_{\mathfrak{p}} = (I R_{\mathfrak{p}})(J R_{\mathfrak{p}})$; and
  $I \neq 0$ implies $I R_{\mathfrak{p}} \neq 0$ (a domain injects into its
  localizations).

  (i) $\Rightarrow$ (v): from $IJ = xR$, extension gives
  $(I R_{\mathfrak{m}})(J R_{\mathfrak{m}}) = x R_{\mathfrak{m}}$ with
  $x \neq 0$.
  (v) $\Rightarrow$ (iii): $R_{\mathfrak{m}}$ is a local domain, so
  Proposition~\ref{infra-invertibility-arith}(g) applies.
  (i) $\Rightarrow$ (ii): the same two steps verbatim at an arbitrary prime
  $\mathfrak{p}$.
  (ii) $\Rightarrow$ (iii) $\Rightarrow$ (iv): specialization.
  (iv) $\Rightarrow$ (iii): if $I \not\subseteq \mathfrak{m}$, any
  $s \in I \setminus \mathfrak{m}$ is a unit of $R_{\mathfrak{m}}$ lying in
  $I R_{\mathfrak{m}}$, so $I R_{\mathfrak{m}} = R_{\mathfrak{m}}$ is principal.

  (iii) $\Rightarrow$ (i): here $I = (x_1, \dots, x_r) \neq 0$. The set
  $I \, I^{-1}$ of finite sums $\sum_k \alpha_k \beta_k$
  ($\alpha_k \in I$, $\beta_k \in I^{-1}$) is an ideal of $R$ by definition of
  $I^{-1}$. Suppose $I \, I^{-1} \neq R$; choose a maximal ideal
  $\mathfrak{m} \supseteq I \, I^{-1}$. By (iii) and clearing denominators,
  $I R_{\mathfrak{m}} = a R_{\mathfrak{m}}$ for some $a \in I$, and $a \neq 0$
  since $I R_{\mathfrak{m}} \neq 0$. For each generator
  $x_j \in I R_{\mathfrak{m}}$ write $x_j = a r_j / s_j$ with $r_j \in R$,
  $s_j \notin \mathfrak{m}$, i.e.\ $s_j x_j \in aR$. With
  $s := s_1 \cdots s_r \notin \mathfrak{m}$ we get $s I \subseteq a R$, hence
  $s/a \in I^{-1}$ and
  $s = a \cdot (s/a) \in I \, I^{-1} \subseteq \mathfrak{m}$, a contradiction.
  Therefore $I \, I^{-1} = R$. Finally, for any nonzero $a \in I$ put
  $J := a I^{-1}$; then $J \subseteq I \, I^{-1} = R$ is an integral ideal and
  $I J = a (I \, I^{-1}) = a R$, proving (i) and the ``moreover'' clause.
  (Under hypothesis (i), finite generation holds by
  Proposition~\ref{infra-invertibility-arith}(f), so the ``moreover'' clause is
  unconditional for invertible ideals.)
\end{proof}

\begin{lemma}[Triviality away from the index]\label{infra-invertibility-away}
  \uses{infra-invertibility-def, notation, infra-index}
  Let $\mathfrak{a} \subseteq \mathcal{O}_d$ be an ideal of finite index
  $n = [\mathcal{O} : \mathfrak{a}]$ and let $\mathfrak{m}$ be a maximal ideal
  of residue characteristic $q$ with $q \nmid n$. Then
  $\mathfrak{a} \not\subseteq \mathfrak{m}$ and
  $\mathfrak{a} \, \mathcal{O}_{\mathfrak{m}} = \mathcal{O}_{\mathfrak{m}}$.
  In particular:
  \begin{enumerate}
    \item[(a)] if $[\mathcal{O} : \mathfrak{a}] = p^m$ for a rational prime $p$
      and $\mathfrak{q}$ lies over a rational prime $q \neq p$, then
      $\mathfrak{a} \, \mathcal{O}_{\mathfrak{q}} = \mathcal{O}_{\mathfrak{q}}$
      is the unit ideal;
    \item[(b)] every maximal ideal containing $\mathfrak{a}$ has residue
      characteristic dividing $n$; for $n = p^m$, $m \geq 1$, every maximal
      ideal containing $\mathfrak{a}$ lies over $p$.
  \end{enumerate}
\end{lemma}

\begin{proof}
  \uses{infra-index}
  By the index layer (Lagrange), $n \in n \mathcal{O} \subseteq \mathfrak{a}$.
  If $n \in \mathfrak{m}$ then $n \in \mathfrak{m} \cap \mathbb{Z} = q\mathbb{Z}$,
  i.e.\ $q \mid n$, contrary to hypothesis. So
  $n \in \mathfrak{a} \setminus \mathfrak{m}$ is a unit of the local ring
  $\mathcal{O}_{\mathfrak{m}}$ lying in
  $\mathfrak{a}\,\mathcal{O}_{\mathfrak{m}}$, whence
  $\mathfrak{a}\,\mathcal{O}_{\mathfrak{m}} = \mathcal{O}_{\mathfrak{m}}$.
  (a) is the case $n = p^m$, $q \neq p$; (b) is the contrapositive, with
  $q \mid p^m$ forcing $q = p$ when $m \geq 1$.
\end{proof}

\begin{corollary}[Invertibility of a $p$-power-index ideal is a condition at $p$]
  \label{infra-invertibility-local-at-p}
  \uses{infra-invertibility, infra-invertibility-away, infra-index}
  Let $p$ be a rational prime and let $\mathfrak{a} \subseteq \mathcal{O}_d$
  have index $p^m$, $m \geq 0$. Then
  $\mathfrak{a}$ is invertible if and only if
  $\mathfrak{a}\,\mathcal{O}_{\mathfrak{p}}$ is principal for every maximal
  ideal $\mathfrak{p}$ of $\mathcal{O}$ lying over $p$.
\end{corollary}

\begin{proof}
  \uses{infra-invertibility, infra-invertibility-away, infra-index}
  Finite index gives $\mathfrak{a} \neq 0$, and $\mathcal{O}_d$ is Noetherian
  (free of rank $2$ over $\mathbb{Z}$), so
  Theorem~\ref{infra-invertibility} applies.
  ($\Rightarrow$) is (i) $\Rightarrow$ (iii) restricted to the primes over $p$.
  ($\Leftarrow$) verifies (iv): a maximal ideal containing $\mathfrak{a}$ lies
  over $p$ by Lemma~\ref{infra-invertibility-away}(b) (for $m = 0$ there is
  none), so its localization is principal by hypothesis; now apply
  (iv) $\Rightarrow$ (i).
\end{proof}

\begin{lemma}[The conductor ideal]\label{infra-invertibility-conductor}
  \uses{notation}
  $f O_K \subseteq \mathcal{O}_d$, and $f O_K$ is simultaneously an ideal of
  $O_K$ and of $\mathcal{O}_d$; moreover $\mathcal{O}_d = \mathbb{Z} + f O_K$.
\end{lemma}

\begin{proof}
  \uses{notation}
  One computes
  $\tau - f\omega = \tfrac{Df^2 + f\sqrt D}{2} - \tfrac{fD + f\sqrt D}{2}
    = \tfrac{Df(f-1)}{2} \in \mathbb{Z}$
  (as $f(f-1)$ is even), so $f\omega \in \mathcal{O}$ and
  $f O_K = f\mathbb{Z} + f\omega\,\mathbb{Z} \subseteq \mathbb{Z} + \tau\mathbb{Z} = \mathcal{O}$
  (also $\tau \in O_K$, so $\mathcal{O} \subseteq O_K$). Being an $O_K$-ideal,
  $f O_K$ is a fortiori closed under multiplication by
  $\mathcal{O} \subseteq O_K$, hence an $\mathcal{O}$-ideal. Finally
  $\mathcal{O} = \mathbb{Z} + \tau \mathbb{Z} \subseteq \mathbb{Z} + f O_K
    \subseteq \mathcal{O}$.
  (Cf.\ Cox, Lemma 7.2.)
\end{proof}

\begin{lemma}[DVR away from the conductor]\label{infra-invertibility-dvr}
  \uses{notation, infra-index, infra-invertibility-conductor}
  Let $\mathfrak{m}$ be a maximal ideal of $\mathcal{O}_d$ of residue
  characteristic $q$ with $q \nmid f$. Then, inside $K$:
  \begin{enumerate}
    \item[(1)] $O_K \subseteq \mathcal{O}_{\mathfrak{m}}$;
    \item[(2)] $\mathfrak{M} := \mathfrak{m}\mathcal{O}_{\mathfrak{m}} \cap O_K$
      is a nonzero prime ideal of $O_K$ with
      $\mathfrak{M} \cap \mathcal{O} = \mathfrak{m}$;
    \item[(3)] $\mathcal{O}_{\mathfrak{m}} = (O_K)_{\mathfrak{M}}$;
    \item[(4)] $\mathcal{O}_{\mathfrak{m}}$ is a discrete valuation ring; in
      particular $\mathfrak{a}\,\mathcal{O}_{\mathfrak{m}}$ is principal for
      every nonzero ideal $\mathfrak{a} \subseteq \mathcal{O}$.
  \end{enumerate}
\end{lemma}

\begin{proof}
  \uses{infra-invertibility-conductor, infra-index}
  (1) $f \notin \mathfrak{m}$, else
  $f \in \mathfrak{m} \cap \mathbb{Z} = q\mathbb{Z}$; so $f$ is a unit of
  $\mathcal{O}_{\mathfrak{m}}$. For $\beta \in O_K$,
  $f\beta \in f O_K \subseteq \mathcal{O}$
  (Lemma~\ref{infra-invertibility-conductor}), hence
  $\beta = (f\beta)/f \in \mathcal{O}_{\mathfrak{m}}$.

  (2) $\mathfrak{m}\mathcal{O}_{\mathfrak{m}}$ is the maximal ideal of the
  local ring $\mathcal{O}_{\mathfrak{m}}$; its contraction $\mathfrak{M}$ to
  the subring $O_K$ is prime, and
  $\mathfrak{M} \cap \mathcal{O}
    = \mathfrak{m}\mathcal{O}_{\mathfrak{m}} \cap \mathcal{O} = \mathfrak{m}$
  by the standard contraction identity. It is nonzero since
  $q \in \mathfrak{m} \cap O_K \subseteq \mathfrak{M}$.

  (3) ($\subseteq$) An element $\alpha/s$ of $\mathcal{O}_{\mathfrak{m}}$
  ($\alpha \in \mathcal{O}$, $s \in \mathcal{O} \setminus \mathfrak{m}$) has
  $s \notin \mathfrak{M}$ — else $s \in \mathfrak{M} \cap \mathcal{O} = \mathfrak{m}$ —
  so $\alpha/s \in (O_K)_{\mathfrak{M}}$.
  ($\supseteq$) An element $\beta/t$ ($\beta \in O_K$,
  $t \in O_K \setminus \mathfrak{M}$) has $\beta, t \in \mathcal{O}_{\mathfrak{m}}$
  by (1) and $t \notin \mathfrak{m}\mathcal{O}_{\mathfrak{m}}$ by definition of
  $\mathfrak{M}$; so $t$ is a unit of $\mathcal{O}_{\mathfrak{m}}$ and
  $\beta/t \in \mathcal{O}_{\mathfrak{m}}$.

  (4) $O_K$ is Dedekind and $\mathfrak{M} \neq 0$ is prime, so
  $(O_K)_{\mathfrak{M}}$ is a DVR (Neukirch I \S 11; Atiyah--Macdonald,
  Thm.\ 9.3). Nonzero ideals of a DVR are powers of the maximal ideal, hence
  principal, and $\mathfrak{a}\,\mathcal{O}_{\mathfrak{m}} \neq 0$.
\end{proof}

\begin{theorem}[Index prime to the conductor implies invertible]
  \label{infra-invertibility-coprime}
  \lean{QuadraticOrder.isInvertibleIdeal_of_coprime_conductor}
  \uses{infra-invertibility-def, notation, infra-index}
  Let $\mathfrak{a} \subseteq \mathcal{O}_d$ be an ideal of finite index
  $n = [\mathcal{O} : \mathfrak{a}]$ with $\gcd(n, f) = 1$. Then $\mathfrak{a}$
  is invertible.
\end{theorem}

\begin{proof}
  \uses{infra-invertibility, infra-invertibility-away, infra-invertibility-dvr, infra-index}
  Finite index gives $\mathfrak{a} \neq 0$, and $\mathcal{O}_d$ is Noetherian;
  we verify condition (iv) of Theorem~\ref{infra-invertibility}. Let
  $\mathfrak{m} \supseteq \mathfrak{a}$ be maximal, of residue characteristic
  $q$. By Lemma~\ref{infra-invertibility-away}(b), $q \mid n$; with
  $\gcd(n, f) = 1$ this forces $q \nmid f$. By
  Lemma~\ref{infra-invertibility-dvr}(4),
  $\mathfrak{a}\,\mathcal{O}_{\mathfrak{m}}$ is principal. Hence (iv) holds and
  $\mathfrak{a}$ is invertible.
\end{proof}

\begin{remark}\label{infra-invertibility-remarks}
  \uses{infra-invertibility, infra-invertibility-coprime,
    infra-invertibility-local-at-p, infra-invertibility-dvr}
  (1) The hypothesis $\gcd([\mathcal{O} : \mathfrak{a}], f) = 1$ is equivalent
  to the classical ``$\mathfrak{a}$ is prime to the conductor'',
  $\mathfrak{a} + f\mathcal{O} = \mathcal{O}$ (Cox, Lemma 7.18): one direction
  is $1 = an + bf \in \mathfrak{a} + f\mathcal{O}$ via
  $n\mathcal{O} \subseteq \mathfrak{a}$; conversely, if a prime $q$ divides
  $\gcd(n, f)$, then $\mathfrak{a} + q\mathcal{O} = \mathcal{O}$ would make
  multiplication by $q$ bijective on $\mathcal{O}/\mathfrak{a}$, contradicting
  $q \mid n$ by Cauchy's theorem; so
  $\mathfrak{a} + f\mathcal{O} \subseteq \mathfrak{a} + q\mathcal{O} \subsetneq \mathcal{O}$.
  (2) The converse of Theorem~\ref{infra-invertibility-coprime} fails: for
  $p \mid f$ the principal ideal $p\mathcal{O}$ has index $p^2$ and is
  invertible. The corrected Corollary~3.1.3* (ERRATA E5) counts exactly the
  invertible ideals of $p$-power index for $p \mid f$, and by
  Corollary~\ref{infra-invertibility-local-at-p} that count is a purely local
  computation at the primes over $p$.
  (3) At a prime over $p \mid f$ the local ring is \emph{not} a DVR, so
  Lemma~\ref{infra-invertibility-dvr} genuinely requires $q \nmid f$; this is
  why the invertible count differs from the total count of Theorem~3.1.2*.
\end{remark}
